% report_qualita_codice.tex — Report di qualità del codice (backend)

\section{Report di Qualità del Codice (Backend)}

\subsection{Strumenti di analisi utilizzati}

L'analisi della qualità del codice è stata condotta utilizzando:
\begin{itemize}
    \item \textbf{Analisi statica manuale} con revisione sistematica del codice
    \item \textbf{SonarLint} integrato nell'IDE per la rilevazione in tempo reale
          di code smell e bug pattern
    \item \textbf{Compilatore Java 23} con warning abilitati
    \item \textbf{JUnit 5 + Mockito} per la verifica della copertura funzionale
\end{itemize}

\emph{Nota: SonarQube Server non è stato configurato per questo progetto.
L'analisi equivalente è stata condotta con gli strumenti sopra descritti.}

\subsection{Metriche generali}

\begin{center}
\begin{tabular}{|l|r|}
\hline
\textbf{Metrica} & \textbf{Valore} \\
\hline
File sorgente Java & 38 \\
\hline
Linee di codice (LOC) & $\sim$2\,000 \\
\hline
Classi & 18 \\
\hline
Interfacce (Repository) & 5 \\
\hline
Enumerazioni & 5 \\
\hline
Record (DTO) & 8 \\
\hline
Package & 7 \\
\hline
Test di unità & 9 (tutti passano) \\
\hline
\end{tabular}
\end{center}

\subsection{Distribuzione del codice per package}

\begin{center}
\small
\begin{tabular}{|l|r|p{7cm}|}
\hline
\textbf{Package} & \textbf{File} & \textbf{Contenuto} \\
\hline
\texttt{model} & 10 & Entità JPA (Bug, User, History, Comment, Label) + 5 enumerazioni \\
\hline
\texttt{service} & 5 & BugService, AuthService, ExportService, AnalyticsService, BugSpecification \\
\hline
\texttt{controller} & 5 & AuthController, BugController, AdminController, ExportController, RestExceptionHandler \\
\hline
\texttt{dto} & 8 & Record Java per request/response \\
\hline
\texttt{repository} & 5 & Interfacce Spring Data JPA \\
\hline
\texttt{config} & 2 & SecurityConfig, DataSeeder \\
\hline
\texttt{security} & 2 & JwtService, JwtAuthFilter \\
\hline
\texttt{root} & 1 & Bugboard26Application (entry point) \\
\hline
\end{tabular}
\end{center}

\subsection{Punti di forza}

\begin{itemize}
    \item \textbf{Struttura dei package chiara}: separazione netta tra model,
    service, controller, dto, repository, config e security, conforme al pattern
    layered architecture.
    \item \textbf{Uso di feature moderne di Java}: record per i DTO (immutabilità,
    riduzione del boilerplate), enumerazioni per i tipi con valori finiti.
    \item \textbf{Best practice Spring Boot}: auto-configuration, starter
    dependencies, gestione centralizzata delle eccezioni tramite
    \texttt{RestExceptionHandler}.
    \item \textbf{Separazione delle responsabilità}: controller ``thin'' che
    delegano immediatamente ai servizi; logica di business concentrata nei
    service.
    \item \textbf{Sicurezza}: autenticazione JWT stateless, password hash con
    BCrypt, configurazione CORS restrittiva.
\end{itemize}

\subsection{Code smell e aree di miglioramento identificate}

\subsubsection{Severità alta}

\begin{enumerate}
    \item \textbf{Autenticazione hardcoded nel controller}: Il metodo
    \texttt{getCurrentUserId()} in \texttt{BugController} restituisce un UUID
    fisso anziché estrarlo dal token JWT. Analogamente,
    \texttt{isCurrentUserAdmin()} ritorna sempre \texttt{true}.
    Questo rappresenta un rischio di sicurezza che deve essere corretto
    prima del deployment.

    \item \textbf{Endpoint \texttt{/me} incompleto}: L'endpoint
    \texttt{GET /api/auth/me} in \texttt{AuthController} non restituisce
    i dati dell'utente corrente in modo completo.
\end{enumerate}

\subsubsection{Severità media}

\begin{enumerate}
    \item \textbf{Metodo lungo}: \texttt{BugService.patch()} ($\sim$41 righe)
    gestisce molteplici responsabilità (validazione permessi, aggiornamento
    campi, gestione duplicati, storico). Potrebbe essere suddiviso in metodi
    più piccoli.

    \item \textbf{Logica duplicata}: la ricerca dell'utente con lancio di
    \texttt{EntityNotFoundException} è ripetuta in più metodi di
    \texttt{BugService}. Estraibile in un metodo helper privato.

    \item \textbf{Query inefficienti}: \texttt{AnalyticsService.getMetrics()}
    e \texttt{getMonthlyReport()} caricano tutti i bug in memoria anziché
    usare query aggregate a livello di database.

    \item \textbf{Possibile NullPointerException}: in \texttt{BugService.patch()},
    l'accesso a \texttt{bug.getAssignee().getId()} non è protetto da null check
    nel caso in cui il bug non abbia un assegnatario.
\end{enumerate}

\subsubsection{Severità bassa}

\begin{enumerate}
    \item \textbf{Magic number}: il limite di 5~MB per gli allegati è espresso
    come \texttt{5 * 1024 * 1024} inline. Dovrebbe essere una costante con nome
    significativo.

    \item \textbf{Codice commentato}: \texttt{BugService.uploadAttachment()}
    contiene righe commentate per il salvataggio su filesystem.

    \item \textbf{Classe utility non finale}: \texttt{BugSpecification} è una
    classe di utility ma non è dichiarata \texttt{final} né ha un costruttore
    privato.
\end{enumerate}

\subsection{Riepilogo}

\begin{center}
\begin{tabular}{|l|c|c|}
\hline
\textbf{Categoria} & \textbf{Conteggio} & \textbf{Severità} \\
\hline
Bug critici (sicurezza) & 2 & Alta \\
\hline
Code smell strutturali & 4 & Media \\
\hline
Code smell minori & 3 & Bassa \\
\hline
\textbf{Totale issue} & \textbf{9} & \\
\hline
\end{tabular}
\end{center}

\textbf{Valutazione complessiva}: il backend presenta una buona organizzazione
architetturale e un uso appropriato dei pattern di Spring Boot. Le issue
critiche sono localizzate e note (autenticazione placeholder), pianificate per
correzione prima del rilascio in produzione. Le issue di severità media e bassa
riguardano refactoring incrementali che non impattano la correttezza funzionale.
